\section{Related Work}\label{s:related}

Previous ABR studies show conflicting algorithm rankings under various conditions. For example, Mao et al. discovered that reinforcement-learning-based Pensieve outperformed traditional ABR algorithms such as MPC-HM, RobustMPC-HM, and BBA \cite{pensieve-sigcomm}, but Yan et al. observed differently in real-world trials \cite{yan2020-situ}. Hybrid schemes have also outperformed pure strategies: Huang et al.’s Stick improved QoE by 44.3\% over linear BBA and 25.9\% over BOLA \cite{stick-infocom20} and outperformed all baselines by 9.5–25.9\% \cite{stick-infocom20}, while DeepMPC similarly outperformed Pensieve by about 4–5.9\% QoE \cite{deepmpc-icip20}. Ghosh et al.’s 5G measurements revealed that under mmWave conditions, RobustMPC was the least stalling ABR (less than 5\%) \cite{5g-variegated}, whereas algorithms like Pensieve and FastMPC experienced increased stalling times by roughly 82\% and 259\% on 5G \cite{5g-variegated}. Pensieve experienced more stalls on 5G while being the best in 4G traces \cite{5g-variegated}, reflecting similarity with our own results. These findings also align with CausalSim’s conclusion that simulator bias can change ABR rankings, cutting error by roughly 50–60\% and producing different insights about ABR algorithms than biased simulators \cite{causalsim-nsdi23}. These comparisons highlight that reported ABR rankings are highly context-dependent. We therefore replay our six ABR algorithms under both Mahimahi and the new CellReplay emulators to see whether prior orderings hold. In our experiments, we also observe shifts like algorithms ranked highest under Mahimahi can drop when run on CellReplay, supporting earlier reports from \cite{5g-variegated}. Overall, ABR performance varies significantly with environment, trace complexity, and evaluation setup. Our results extend this understanding by showing that emulator choice and trace variability can significantly affect ABR rankings. 
Hybrid or rule-based methods (TTP and linear BBA) remain strong and balanced across setups, whereas learned policies like Pensieve must be carefully tuned to each environment.

Comparing our CellReplay‐based ABR rankings to prior work reveals some deviations and similarities. For example, buffer-based BOLA in its original evaluation achieved 84–95\% of the ideal performance~\cite{bola-bola}, but in our Mahimahi-based experiments, BOLA became the weakest algorithm under constrained 4G, and became the smoothest algorithm under CellReplay-based 5G settings. This suggests that network conditions can change ABR rankings. Similarly, SENSEI’s user-sensitivity model shows that weighting quality by viewer importance can reorder outcomes. Their evaluation reports roughly a 15\% QoE gain by adapting to content sensitivity~\cite{sensei-nsdi21}, implying that ABR rankings may change under personalised QoE metrics. Finally, new codec-aware systems like Swift and LiFteR shift adaptation to the coding layer. Swift’s layered neural codec~\cite{swift-nsdi22} , and LiFteR’s loose frame-referencing (parallel learned-codec decoding) achieve much higher frame rates and quality~\cite{lifter-nsdi24}. These innovations suggest future ABRs may choose which layers or frames to fetch, potentially altering relative performance. In summary, while our general trends agree with past studies, using CellReplay and accounting for user preferences or new codec designs can reveal different ABR rankings and merit further investigation.

%%% Local Variables:
%%% mode: latex
%%% TeX-master: "../thesis"
%%% End:
